%% bare_jrnl_compsoc.tex
%% V1.4
%% 2012/12/27
%% by Michael Shell
%% See:
%% http://www.michaelshell.org/
%% for current contact information.
%%
%% This is a skeleton file demonstrating the use of IEEEtran.cls
%% (requires IEEEtran.cls version 1.8 or later) with an IEEE Computer
%% Society journal paper.
%%
%% Support sites:
%% http://www.michaelshell.org/tex/ieeetran/
%% http://www.ctan.org/tex-archive/macros/latex/contrib/IEEEtran/
%% and
%% http://www.ieee.org/

%%*************************************************************************
%% Legal Notice:
%% This code is offered as-is without any warranty either expressed or
%% implied; without even the implied warranty of MERCHANTABILITY or
%% FITNESS FOR A PARTICULAR PURPOSE! 
%% User assumes all risk.
%% In no event shall IEEE or any contributor to this code be liable for
%% any damages or losses, including, but not limited to, incidental,
%% consequential, or any other damages, resulting from the use or misuse
%% of any information contained here.
%%
%% All comments are the opinions of their respective authors and are not
%% necessarily endorsed by the IEEE.
%%
%% This work is distributed under the LaTeX Project Public License (LPPL)
%% ( http://www.latex-project.org/ ) version 1.3, and may be freely used,
%% distributed and modified. A copy of the LPPL, version 1.3, is included
%% in the base LaTeX documentation of all distributions of LaTeX released
%% 2003/12/01 or later.
%% Retain all contribution notices and credits.
%% ** Modified files should be clearly indicated as such, including  **
%% ** renaming them and changing author support contact information. **
%%
%% File list of work: IEEEtran.cls, IEEEtran_HOWTO.pdf, bare_adv.tex,
%%                    bare_conf.tex, bare_jrnl.tex, bare_jrnl_compsoc.tex,
%%                    bare_jrnl_transmag.tex
%%*************************************************************************

% *** Authors should verify (and, if needed, correct) their LaTeX system  ***
% *** with the testflow diagnostic prior to trusting their LaTeX platform ***
% *** with production work. IEEE's font choices can trigger bugs that do  ***
% *** not appear when using other class files.                            ***
% The testflow support page is at:
% http://www.michaelshell.org/tex/testflow/




% Note that the a4paper option is mainly intended so that authors in
% countries using A4 can easily print to A4 and see how their papers will
% look in print - the typesetting of the document will not typically be
% affected with changes in paper size (but the bottom and side margins will).
% Use the testflow package mentioned above to verify correct handling of
% both paper sizes by the user's LaTeX system.
%
% Also note that the "draftcls" or "draftclsnofoot", not "draft", option
% should be used if it is desired that the figures are to be displayed in
% draft mode.
%
% The Computer Society usually requires 12pt for submissions.
%
\documentclass[12pt,journal,compsoc]{IEEEtran}
%
% If IEEEtran.cls has not been installed into the LaTeX system files,
% manually specify the path to it like:
% \documentclass[12pt,journal,compsoc]{../sty/IEEEtran}





% Some very useful LaTeX packages include:
% (uncomment the ones you want to load)


% *** MISC UTILITY PACKAGES ***
%
%\usepackage{ifpdf}
% Heiko Oberdiek's ifpdf.sty is very useful if you need conditional
% compilation based on whether the output is pdf or dvi.
% usage:
% \ifpdf
%   % pdf code
% \else
%   % dvi code
% \fi
% The latest version of ifpdf.sty can be obtained from:
% http://www.ctan.org/tex-archive/macros/latex/contrib/oberdiek/
% Also, note that IEEEtran.cls V1.7 and later provides a builtin
% \ifCLASSINFOpdf conditional that works the same way.
% When switching from latex to pdflatex and vice-versa, the compiler may
% have to be run twice to clear warning/error messages.






% *** CITATION PACKAGES ***
%
%\ifCLASSOPTIONcompsoc
  % IEEE Computer Society needs nocompress option
  % requires cite.sty v4.0 or later (November 2003)
  %\usepackage[nocompress]{cite}
%\else
  % normal IEEE
  %\usepackage{cite}
%\fi
% cite.sty was written by Donald Arseneau
% V1.6 and later of IEEEtran pre-defines the format of the cite.sty package
% \cite{} output to follow that of IEEE. Loading the cite package will
% result in citation numbers being automatically sorted and properly
% "compressed/ranged". e.g., [1], [9], [2], [7], [5], [6] without using
% cite.sty will become [1], [2], [5]--[7], [9] using cite.sty. cite.sty's
% \cite will automatically add leading space, if needed. Use cite.sty's
% noadjust option (cite.sty V3.8 and later) if you want to turn this off
% such as if a citation ever needs to be enclosed in parenthesis.
% cite.sty is already installed on most LaTeX systems. Be sure and use
% version 4.0 (2003-05-27) and later if using hyperref.sty. cite.sty does
% not currently provide for hyperlinked citations.
% The latest version can be obtained at:
% http://www.ctan.org/tex-archive/macros/latex/contrib/cite/
% The documentation is contained in the cite.sty file itself.
%
% Note that some packages require special options to format as the Computer
% Society requires. In particular, Computer Society  papers do not use
% compressed citation ranges as is done in typical IEEE papers
% (e.g., [1]-[4]). Instead, they list every citation separately in order
% (e.g., [1], [2], [3], [4]). To get the latter we need to load the cite
% package with the nocompress option which is supported by cite.sty v4.0
% and later. Note also the use of a CLASSOPTION conditional provided by
% IEEEtran.cls V1.7 and later.





% *** GRAPHICS RELATED PACKAGES ***
%
\ifCLASSINFOpdf
  \usepackage[pdftex]{graphicx}
  % declare the path(s) where your graphic files are
  % \graphicspath{{../pdf/}{../jpeg/}}
  \graphicspath{{img/}}
  % and their extensions so you won't have to specify these with
  % every instance of \includegraphics
  \DeclareGraphicsExtensions{.pdf,.jpeg,.png,.eps}
\else
  % or other class option (dvipsone, dvipdf, if not using dvips). graphicx
  % will default to the driver specified in the system graphics.cfg if no
  % driver is specified.
  \usepackage[dvips]{graphicx}
  % declare the path(s) where your graphic files are
  \graphicspath{{img/}}
  % and their extensions so you won't have to specify these with
  % every instance of \includegraphics
  \DeclareGraphicsExtensions{.eps,.jpeg,.png,.pdf}
\fi
% graphicx was written by David Carlisle and Sebastian Rahtz. It is
% required if you want graphics, photos, etc. graphicx.sty is already
% installed on most LaTeX systems. The latest version and documentation
% can be obtained at: 
% http://www.ctan.org/tex-archive/macros/latex/required/graphics/
% Another good source of documentation is "Using Imported Graphics in
% LaTeX2e" by Keith Reckdahl which can be found at:
% http://www.ctan.org/tex-archive/info/epslatex/
%
% latex, and pdflatex in dvi mode, support graphics in encapsulated
% postscript (.eps) format. pdflatex in pdf mode supports graphics
% in .pdf, .jpeg, .png and .mps (metapost) formats. Users should ensure
% that all non-photo figures use a vector format (.eps, .pdf, .mps) and
% not a bitmapped formats (.jpeg, .png). IEEE frowns on bitmapped formats
% which can result in "jaggedy"/blurry rendering of lines and letters as
% well as large increases in file sizes.
%
% You can find documentation about the pdfTeX application at:
% http://www.tug.org/applications/pdftex






% *** MATH PACKAGES ***
%
\usepackage[cmex10]{amsmath}
% A popular package from the American Mathematical Society that provides
% many useful and powerful commands for dealing with mathematics. If using
% it, be sure to load this package with the cmex10 option to ensure that
% only type 1 fonts will utilized at all point sizes. Without this option,
% it is possible that some math symbols, particularly those within
% footnotes, will be rendered in bitmap form which will result in a
% document that can not be IEEE Xplore compliant!
%
% Also, note that the amsmath package sets \interdisplaylinepenalty to 10000
% thus preventing page breaks from occurring within multiline equations. Use:
\interdisplaylinepenalty=2500
% after loading amsmath to restore such page breaks as IEEEtran.cls normally
% does. amsmath.sty is already installed on most LaTeX systems. The latest
% version and documentation can be obtained at:
% http://www.ctan.org/tex-archive/macros/latex/required/amslatex/math/





% *** SPECIALIZED LIST PACKAGES ***
%
%\usepackage{algorithmic}
% algorithmic.sty was written by Peter Williams and Rogerio Brito.
% This package provides an algorithmic environment fo describing algorithms.
% You can use the algorithmic environment in-text or within a figure
% environment to provide for a floating algorithm. Do NOT use the algorithm
% floating environment provided by algorithm.sty (by the same authors) or
% algorithm2e.sty (by Christophe Fiorio) as IEEE does not use dedicated
% algorithm float types and packages that provide these will not provide
% correct IEEE style captions. The latest version and documentation of
% algorithmic.sty can be obtained at:
% http://www.ctan.org/tex-archive/macros/latex/contrib/algorithms/
% There is also a support site at:
% http://algorithms.berlios.de/index.html
% Also of interest may be the (relatively newer and more customizable)
% algorithmicx.sty package by Szasz Janos:
% http://www.ctan.org/tex-archive/macros/latex/contrib/algorithmicx/




% *** ALIGNMENT PACKAGES ***
%
\usepackage{array}
% Frank Mittelbach's and David Carlisle's array.sty patches and improves
% the standard LaTeX2e array and tabular environments to provide better
% appearance and additional user controls. As the default LaTeX2e table
% generation code is lacking to the point of almost being broken with
% respect to the quality of the end results, all users are strongly
% advised to use an enhanced (at the very least that provided by array.sty)
% set of table tools. array.sty is already installed on most systems. The
% latest version and documentation can be obtained at:
% http://www.ctan.org/tex-archive/macros/latex/required/tools/


% IEEEtran contains the IEEEeqnarray family of commands that can be used to
% generate multiline equations as well as matrices, tables, etc., of high
% quality.




% *** SUBFIGURE PACKAGES ***
\ifCLASSOPTIONcompsoc
  \usepackage[caption=false,font=normalsize,labelfont=sf,textfont=sf]{subfig}
\else
  \usepackage[caption=false,font=footnotesize]{subfig}
\fi
% subfig.sty, written by Steven Douglas Cochran, is the modern replacement
% for subfigure.sty, the latter of which is no longer maintained and is
% incompatible with some LaTeX packages including fixltx2e. However,
% subfig.sty requires and automatically loads Axel Sommerfeldt's caption.sty
% which will override IEEEtran.cls' handling of captions and this will result
% in non-IEEE style figure/table captions. To prevent this problem, be sure
% and invoke subfig.sty's "caption=false" package option (available since
% subfig.sty version 1.3, 2005/06/28) as this is will preserve IEEEtran.cls
% handling of captions.
% Note that the Computer Society format requires a larger sans serif font
% than the serif footnote size font used in traditional IEEE formatting
% and thus the need to invoke different subfig.sty package options depending
% on whether compsoc mode has been enabled.
%
% The latest version and documentation of subfig.sty can be obtained at:
% http://www.ctan.org/tex-archive/macros/latex/contrib/subfig/




% *** FLOAT PACKAGES ***
%
\usepackage{fixltx2e}
% fixltx2e, the successor to the earlier fix2col.sty, was written by
% Frank Mittelbach and David Carlisle. This package corrects a few problems
% in the LaTeX2e kernel, the most notable of which is that in current
% LaTeX2e releases, the ordering of single and double column floats is not
% guaranteed to be preserved. Thus, an unpatched LaTeX2e can allow a
% single column figure to be placed prior to an earlier double column
% figure. The latest version and documentation can be found at:
% http://www.ctan.org/tex-archive/macros/latex/base/


%\usepackage{stfloats}
% stfloats.sty was written by Sigitas Tolusis. This package gives LaTeX2e
% the ability to do double column floats at the bottom of the page as well
% as the top. (e.g., "\begin{figure*}[!b]" is not normally possible in
% LaTeX2e). It also provides a command:
%\fnbelowfloat
% to enable the placement of footnotes below bottom floats (the standard
% LaTeX2e kernel puts them above bottom floats). This is an invasive package
% which rewrites many portions of the LaTeX2e float routines. It may not work
% with other packages that modify the LaTeX2e float routines. The latest
% version and documentation can be obtained at:
% http://www.ctan.org/tex-archive/macros/latex/contrib/sttools/
% Do not use the stfloats baselinefloat ability as IEEE does not allow
% \baselineskip to stretch. Authors submitting work to the IEEE should note
% that IEEE rarely uses double column equations and that authors should try
% to avoid such use. Do not be tempted to use the cuted.sty or midfloat.sty
% packages (also by Sigitas Tolusis) as IEEE does not format its papers in
% such ways.
% Do not attempt to use stfloats with fixltx2e as they are incompatible.
% Instead, use Morten Hogholm'a dblfloatfix which combines the features
% of both fixltx2e and stfloats:
%
\usepackage{dblfloatfix}
% The latest version can be found at:
% http://www.ctan.org/tex-archive/macros/latex/contrib/dblfloatfix/




%\ifCLASSOPTIONcaptionsoff
%  \usepackage[nomarkers]{endfloat}
% \let\MYoriglatexcaption\caption
% \renewcommand{\caption}[2][\relax]{\MYoriglatexcaption[#2]{#2}}
%\fi
% endfloat.sty was written by James Darrell McCauley, Jeff Goldberg and 
% Axel Sommerfeldt. This package may be useful when used in conjunction with 
% IEEEtran.cls'  captionsoff option. Some IEEE journals/societies require that
% submissions have lists of figures/tables at the end of the paper and that
% figures/tables without any captions are placed on a page by themselves at
% the end of the document. If needed, the draftcls IEEEtran class option or
% \CLASSINPUTbaselinestretch interface can be used to increase the line
% spacing as well. Be sure and use the nomarkers option of endfloat to
% prevent endfloat from "marking" where the figures would have been placed
% in the text. The two hack lines of code above are a slight modification of
% that suggested by in the endfloat docs (section 8.4.1) to ensure that
% the full captions always appear in the list of figures/tables - even if
% the user used the short optional argument of \caption[]{}.
% IEEE papers do not typically make use of \caption[]'s optional argument,
% so this should not be an issue. A similar trick can be used to disable
% captions of packages such as subfig.sty that lack options to turn off
% the subcaptions:
% For subfig.sty:
% \let\MYorigsubfloat\subfloat
% \renewcommand{\subfloat}[2][\relax]{\MYorigsubfloat[]{#2}}
% However, the above trick will not work if both optional arguments of
% the \subfloat command are used. Furthermore, there needs to be a
% description of each subfigure *somewhere* and endfloat does not add
% subfigure captions to its list of figures. Thus, the best approach is to
% avoid the use of subfigure captions (many IEEE journals avoid them anyway)
% and instead reference/explain all the subfigures within the main caption.
% The latest version of endfloat.sty and its documentation can obtained at:
% http://www.ctan.org/tex-archive/macros/latex/contrib/endfloat/
%
% The IEEEtran \ifCLASSOPTIONcaptionsoff conditional can also be used
% later in the document, say, to conditionally put the References on a 
% page by themselves.




% *** PDF, URL AND HYPERLINK PACKAGES ***
%
%\usepackage{url}
% url.sty was written by Donald Arseneau. It provides better support for
% handling and breaking URLs. url.sty is already installed on most LaTeX
% systems. The latest version and documentation can be obtained at:
% http://www.ctan.org/tex-archive/macros/latex/contrib/url/
% Basically, \url{my_url_here}.





% *** Do not adjust lengths that control margins, column widths, etc. ***
% *** Do not use packages that alter fonts (such as pslatex).         ***
% There should be no need to do such things with IEEEtran.cls V1.6 and later.
% (Unless specifically asked to do so by the journal or conference you plan
% to submit to, of course. )


% correct bad hyphenation here
\hyphenation{op-tical net-works semi-conduc-tor}


%SACHA
%Paquete para incluir links (ya sea para mandar un mail o para p\'aginas web)
\usepackage{hyperref}
%\usepackage{nameref}

%http://en.wikibooks.org/wiki/LaTeX/Importing_Graphics
\usepackage{pdfpages}

%este paquete permite incluir acentos. Notar que espera un formato ANSI-blah de archivo. 
%Si en lugar de eso se tiene un utf8 (usual en los linux),
%entonces usar \usepackage[utf8]{inputenc}
\usepackage[utf8]{inputenc}
\usepackage[latin,spanish]{babel}
\usepackage{csquotes}

%http://en.wikibooks.org/wiki/LaTeX/Floats,_Figures_and_Captions#Keeping_floats_in_their_place
\usepackage{float}
%\usepackage{morefloats}
\usepackage{placeins}
%\usepackage{cuted}
%\usepackage{capt-of}

%Para manipular imagenes y muchas cosas con Tikz
%\usepackage{pgf}
%\usepackage{tikz}
%\usetikzlibrary{arrows,automata}
%\usetikzlibrary{spy} %paquete para poder hacer zoom

%http://tex.stackexchange.com/questions/25414/how-to-create-magnified-subfigures-and-corresponding-boxes-for-portions-of-a-lar
%\usepackage{graphicx}
\usepackage{caption}
\usepackage{subcaption}
\usepackage{tikz}
\usepackage{pgfplots}
\usetikzlibrary{spy,calc}
\usepackage{hyperref}


\newif\ifblackandwhitecycle
\gdef\patternnumber{0}

\pgfkeys{/tikz/.cd,
    zoombox paths/.style={
        draw=orange,
        very thick
    },
    black and white/.is choice,
    black and white/.default=static,
    black and white/static/.style={ 
        draw=white,   
        zoombox paths/.append style={
            draw=white,
            postaction={
                draw=black,
                loosely dashed
            }
        }
    },
    black and white/static/.code={
        \gdef\patternnumber{1}
    },
    black and white/cycle/.code={
        \blackandwhitecycletrue
        \gdef\patternnumber{1}
    },
    black and white pattern/.is choice,
    black and white pattern/0/.style={},
    black and white pattern/1/.style={    
            draw=white,
            postaction={
                draw=black,
                dash pattern=on 2pt off 2pt
            }
    },
    black and white pattern/2/.style={    
            draw=white,
            postaction={
                draw=black,
                dash pattern=on 4pt off 4pt
            }
    },
    black and white pattern/3/.style={    
            draw=white,
            postaction={
                draw=black,
                dash pattern=on 4pt off 4pt on 1pt off 4pt
            }
    },
    black and white pattern/4/.style={    
            draw=white,
            postaction={
                draw=black,
                dash pattern=on 4pt off 2pt on 2 pt off 2pt on 2 pt off 2pt
            }
    },
    zoomboxarray inner gap/.initial=5pt,
    zoomboxarray columns/.initial=2,
    zoomboxarray rows/.initial=2,
    subfigurename/.initial={},
    figurename/.initial={zoombox},
    zoomboxarray/.style={
        execute at begin picture={
            \begin{scope}[
                spy using outlines={%
                    zoombox paths,
                    width=\imagewidth / \pgfkeysvalueof{/tikz/zoomboxarray columns} - (\pgfkeysvalueof{/tikz/zoomboxarray columns} - 1) / \pgfkeysvalueof{/tikz/zoomboxarray columns} * \pgfkeysvalueof{/tikz/zoomboxarray inner gap} -\pgflinewidth,
                    height=\imageheight / \pgfkeysvalueof{/tikz/zoomboxarray rows} - (\pgfkeysvalueof{/tikz/zoomboxarray rows} - 1) / \pgfkeysvalueof{/tikz/zoomboxarray rows} * \pgfkeysvalueof{/tikz/zoomboxarray inner gap}-\pgflinewidth,
                    magnification=3,
                    every spy on node/.style={
                        zoombox paths
                    },
                    every spy in node/.style={
                        zoombox paths
                    }
                }
            ]
        },
        execute at end picture={
            \end{scope}
            \node at (image.north) [anchor=north,inner sep=0pt] {\subcaptionbox{\label{\pgfkeysvalueof{/tikz/figurename}-image}}{\phantomimage}};
            \node at (zoomboxes container.north) [anchor=north,inner sep=0pt] {\subcaptionbox{\label{\pgfkeysvalueof{/tikz/figurename}-zoom}}{\phantomimage}};
     \gdef\patternnumber{0}
        },
        spymargin/.initial=0.5em,
        zoomboxes xshift/.initial=1,
        zoomboxes right/.code=\pgfkeys{/tikz/zoomboxes xshift=1},
        zoomboxes left/.code=\pgfkeys{/tikz/zoomboxes xshift=-1},
        zoomboxes yshift/.initial=0,
        zoomboxes above/.code={
            \pgfkeys{/tikz/zoomboxes yshift=1},
            \pgfkeys{/tikz/zoomboxes xshift=0}
        },
        zoomboxes below/.code={
            \pgfkeys{/tikz/zoomboxes yshift=-1},
            \pgfkeys{/tikz/zoomboxes xshift=0}
        },
        caption margin/.initial=4ex,
    },
    adjust caption spacing/.code={},
    image container/.style={
        inner sep=0pt,
        at=(image.north),
        anchor=north,
        adjust caption spacing
    },
    zoomboxes container/.style={
        inner sep=0pt,
        at=(image.north),
        anchor=north,
        name=zoomboxes container,
        xshift=\pgfkeysvalueof{/tikz/zoomboxes xshift}*(\imagewidth+\pgfkeysvalueof{/tikz/spymargin}),
        yshift=\pgfkeysvalueof{/tikz/zoomboxes yshift}*(\imageheight+\pgfkeysvalueof{/tikz/spymargin}+\pgfkeysvalueof{/tikz/caption margin}),
        adjust caption spacing
    },
    calculate dimensions/.code={
        \pgfpointdiff{\pgfpointanchor{image}{south west} }{\pgfpointanchor{image}{north east} }
        \pgfgetlastxy{\imagewidth}{\imageheight}
        \global\let\imagewidth=\imagewidth
        \global\let\imageheight=\imageheight
        \gdef\columncount{1}
        \gdef\rowcount{1}
        \gdef\zoomboxcount{1}
    },
    image node/.style={
        inner sep=0pt,
        name=image,
        anchor=south west,
        append after command={
            [calculate dimensions]
            node [image container,subfigurename=\pgfkeysvalueof{/tikz/figurename}-image] {\phantomimage}
            node [zoomboxes container,subfigurename=\pgfkeysvalueof{/tikz/figurename}-zoom] {\phantomimage}
        }
    },
    color code/.style={
        zoombox paths/.append style={draw=#1}
    },
    connect zoomboxes/.style={
    spy connection path={\draw[draw=none,zoombox paths] (tikzspyonnode) -- (tikzspyinnode);}
    },
    help grid code/.code={
        \begin{scope}[
                x={(image.south east)},
                y={(image.north west)},
                font=\footnotesize,
                help lines,
                overlay
            ]
            \foreach \x in {0,1,...,9} { 
                \draw(\x/10,0) -- (\x/10,1);
                \node [anchor=north] at (\x/10,0) {0.\x};
            }
            \foreach \y in {0,1,...,9} {
                \draw(0,\y/10) -- (1,\y/10);                        \node [anchor=east] at (0,\y/10) {0.\y};
            }
        \end{scope}    
    },
    help grid/.style={
        append after command={
            [help grid code]
        }
    },
}

\newcommand\phantomimage{%
    \phantom{%
        \rule{\imagewidth}{\imageheight}%
    }%
}
\newcommand\zoombox[2][]{
    \begin{scope}[zoombox paths]
        \pgfmathsetmacro\xpos{
            (\columncount-1)*(\imagewidth / \pgfkeysvalueof{/tikz/zoomboxarray columns} + \pgfkeysvalueof{/tikz/zoomboxarray inner gap} / \pgfkeysvalueof{/tikz/zoomboxarray columns} ) + \pgflinewidth
        }
        \pgfmathsetmacro\ypos{
            (\rowcount-1)*( \imageheight / \pgfkeysvalueof{/tikz/zoomboxarray rows} + \pgfkeysvalueof{/tikz/zoomboxarray inner gap} / \pgfkeysvalueof{/tikz/zoomboxarray rows} ) + 0.5*\pgflinewidth
        }
        \edef\dospy{\noexpand\spy [
            #1,
            zoombox paths/.append style={
                black and white pattern=\patternnumber
            },
            every spy on node/.append style={#1},
            x=\imagewidth,
            y=\imageheight
        ] on (#2) in node [anchor=north west] at ($(zoomboxes container.north west)+(\xpos pt,-\ypos pt)$);}
        \dospy
        \pgfmathtruncatemacro\pgfmathresult{ifthenelse(\columncount==\pgfkeysvalueof{/tikz/zoomboxarray columns},\rowcount+1,\rowcount)}
        \global\let\rowcount=\pgfmathresult
        \pgfmathtruncatemacro\pgfmathresult{ifthenelse(\columncount==\pgfkeysvalueof{/tikz/zoomboxarray columns},1,\columncount+1)}
        \global\let\columncount=\pgfmathresult
        \ifblackandwhitecycle
            \pgfmathtruncatemacro{\newpatternnumber}{\patternnumber+1}
            \global\edef\patternnumber{\newpatternnumber}
        \fi
    \end{scope}
}

%	\begin{figure}[H]
%		\begin{tikzpicture}[zoomboxarray,
%		    zoomboxes below,
%		    zoomboxarray columns=1,
%		    zoomboxarray rows=1,
%		    connect zoomboxes,
%		    zoombox paths/.append style={ultra thick, red}]
%		    \node [image node] { \includegraphics[width=0.5\textwidth]{img/graph/graph} };
%		    \zoombox[magnification=20]{0.38,0.5}
%		\end{tikzpicture}
%	\end{figure}

%http://tex.stackexchange.com/questions/29003/pagebreak-and-two-column-text?rq=1
%\usepackage{flushend}

%http://en.wikibooks.org/wiki/LaTeX/List_Structures#Inline_lists
\usepackage[inline]{enumitem}

%Bibliografia
\usepackage[backend=biber,style=numeric-comp]{biblatex}
\bibliography{ieeecls.bib}

%Footnotes
%http://en.wikibooks.org/wiki/LaTeX/Footnotes_and_Margin_Notes
\renewcommand{\thefootnote}{\alph{footnote}}
\usepackage{perpage} %the perpage package
\MakePerPage{footnote} %the perpage package command


%------------------------------------------------------------------------

\begin{document}
%\includepdf{Caratula_Redes_TP1.pdf}
%
% paper title
% can use linebreaks \\ within to get better formatting as desired
% Do not put math or special symbols in the title.
%\title{Bare Demo of IEEEtran.cls\\ for Computer Society Journals}
\title{Trabajo Pr\'actico 3: PTC}
%
%
% author names and IEEE memberships
% note positions of commas and nonbreaking spaces ( ~ ) LaTeX will not break
% a structure at a ~ so this keeps an author's name from being broken across
% two lines.
% use \thanks{} to gain access to the first footnote area
% a separate \thanks must be used for each paragraph as LaTeX2e's \thanks
% was not built to handle multiple paragraphs
%
%
%\IEEEcompsocitemizethanks is a special \thanks that produces the bulleted
% lists the Computer Society journals use for "first footnote" author
% affiliations. Use \IEEEcompsocthanksitem which works much like \item
% for each affiliation group. When not in compsoc mode,
% \IEEEcompsocitemizethanks becomes like \thanks and
% \IEEEcompsocthanksitem becomes a line break with idention. This
% facilitates dual compilation, although admittedly the differences in the
% desired content of \author between the different types of papers makes a
% one-size-fits-all approach a daunting prospect. For instance, compsoc 
% journal papers have the author affiliations above the "Manuscript
% received ..."  text while in non-compsoc journals this is reversed. Sigh.

%\author{Michael~Shell,~\IEEEmembership{Member,~IEEE,}
%        John~Doe,~\IEEEmembership{Fellow,~OSA,}
%        and~Jane~Doe,~\IEEEmembership{Life~Fellow,~IEEE}% <-this % stops a space
%\IEEEcompsocitemizethanks{\IEEEcompsocthanksitem M. Shell is with the Department
%of Electrical and Computer Engineering, Georgia Institute of Technology, Atlanta,
%GA, 30332.\protect\\
% note need leading \protect in front of \\ to get a newline within \thanks as
% \\ is fragile and will error, could use \hfil\break instead.
%E-mail: see http://www.michaelshell.org/contact.html
%\IEEEcompsocthanksitem J. Doe and J. Doe are with Anonymous University.}% <-this % stops an unwanted space
%\thanks{Manuscript received April 19, 2005; revised December 27, 2012.}}

% note the % following the last \IEEEmembership and also \thanks - 
% these prevent an unwanted space from occurring between the last author name
% and the end of the author line. i.e., if you had this:
% 
% \author{....lastname \thanks{...} \thanks{...} }
%                     ^------------^------------^----Do not want these spaces!
%
% a space would be appended to the last name and could cause every name on that
% line to be shifted left slightly. This is one of those "LaTeX things". For
% instance, "\textbf{A} \textbf{B}" will typeset as "A B" not "AB". To get
% "AB" then you have to do: "\textbf{A}\textbf{B}"
% \thanks is no different in this regard, so shield the last } of each \thanks
% that ends a line with a % and do not let a space in before the next \thanks.
% Spaces after \IEEEmembership other than the last one are OK (and needed) as
% you are supposed to have spaces between the names. For what it is worth,
% this is a minor point as most people would not even notice if the said evil
% space somehow managed to creep in.

\author{%
	Federico E. Allocati,
    Sacha S. Kantor
    y Mat\'ias Schmit
}



% The paper headers
%\markboth{Journal of \LaTeX\ Class Files,~Vol.~11, No.~4, December~2012}%
%{Shell \MakeLowercase{\textit{et al.}}: Bare Demo of IEEEtran.cls for Computer Society Journals}
% The only time the second header will appear is for the odd numbered pages
% after the title page when using the twoside option.
% 
% *** Note that you probably will NOT want to include the author's ***
% *** name in the headers of peer review papers.                   ***
% You can use \ifCLASSOPTIONpeerreview for conditional compilation here if
% you desire.
\markboth{Teor\'ia de las Comunicaciones}{, Departamento de Computaci\'on, FCEyN - UBA, Noviembre 2014}



% The publisher's ID mark at the bottom of the page is less important with
% Computer Society journal papers as those publications place the marks
% outside of the main text columns and, therefore, unlike regular IEEE
% journals, the available text space is not reduced by their presence.
% If you want to put a publisher's ID mark on the page you can do it like
% this:
%\IEEEpubid{0000--0000/00\$00.00~\copyright~2012 IEEE}
% or like this to get the Computer Society new two part style.
%\IEEEpubid{\makebox[\columnwidth]{\hfill 0000--0000/00/\$00.00~\copyright~2012 IEEE}%
%\hspace{\columnsep}\makebox[\columnwidth]{Published by the IEEE Computer Society\hfill}}
% Remember, if you use this you must call \IEEEpubidadjcol in the second
% column for its text to clear the IEEEpubid mark (Computer Society jorunal
% papers don't need this extra clearance.)



% use for special paper notices
%\IEEEspecialpapernotice{(Invited Paper)}




% for Computer Society papers, we must declare the abstract and index terms
% PRIOR to the title within the \IEEEtitleabstractindextext IEEEtran
% command as these need to go into the title area created by \maketitle.
% As a general rule, do not put math, special symbols or citations
% in the abstract or keywords.
\IEEEtitleabstractindextext{%
\begin{abstract}
    Esta segunda parte de la materia nos propone experimentar con ciertos
aspectos ya en un nivel m\'as "alto" que la del trabajo anterior. Mientras
que antes nos ocupamos de estudiar los comportamientos de las redes locales
(para un \textit{scope} particular, que era el de los paquetes \textit{ARP}),
ahora nos concentramos ya en redes m\'as "grandes", donde se trabaja
con el protocolo \textit{IP} y ruteo de paquetes, y otro protocolo utilizado
para poder diagnosticar redes desde el punto de vista de la conectividad y el
ruteo: \textit{ICMP}.

\par As\'i pues, el objetivo de este trabajo es la de mejorar la comprensi\'on
de dichos protocolos y alguno de sus usos, as\'i como tambi\'en poder
estudiar topolog\'ias de ruteo\cite{routing} reales e identificar sus puntos
importantes (para alguna definici\'on de \textit{importante}).

\end{abstract}

% Note that keywords are not normally used for peerreview papers.
\begin{IEEEkeywords}
    Teor\'ia de las Comunicaciones, UBA, FCEyN, TCP, PTC, End-To-End,
    Conexion-Oriented, TP3
\end{IEEEkeywords}}


% make the title area
\maketitle


% To allow for easy dual compilation without having to reenter the
% abstract/keywords data, the \IEEEtitleabstractindextext text will
% not be used in maketitle, but will appear (i.e., to be "transported")
% here as \IEEEdisplaynontitleabstractindextext when the compsoc 
% or transmag modes are not selected <OR> if conference mode is selected 
% - because all conference papers position the abstract like regular
% papers do.
\IEEEdisplaynontitleabstractindextext
% \IEEEdisplaynontitleabstractindextext has no effect when using
% compsoc or transmag under a non-conference mode.



% For peer review papers, you can put extra information on the cover
% page as needed:
% \ifCLASSOPTIONpeerreview
% \begin{center} \bfseries EDICS Category: 3-BBND \end{center}
% \fi
%
% For peerreview papers, this IEEEtran command inserts a page break and
% creates the second title. It will be ignored for other modes.
\IEEEpeerreviewmaketitle



%\section{Introduction}
% Computer Society journal papers do something a tad strange with the very
% first section heading (almost always called "Introduction"). They place it
% ABOVE the main text! IEEEtran.cls currently does not do this for you.
% However, You can achieve this effect by making LaTeX jump through some
% hoops via something like:
%
\ifCLASSOPTIONcompsoc
  \noindent\raisebox{2\baselineskip}[0pt][0pt]%
  {\parbox{\columnwidth}{\section*{Introducci\'on y Aclaraciones}\label{sec:introduccion}%
  \global\everypar=\everypar}}%
  \vspace{-1\baselineskip}\vspace{-\parskip}\par
\else
  \section*{Introducci\'on y Aclaraciones}\label{sec:introduccion}\par
%  \section{Introduction}\label{sec:introduction}\par
\fi
\IEEEPARstart{E}{ste} trabajo se concentra principalmente en el an\'alisis

%
% Admittedly, this is a hack and may well be fragile, but seems to do the
% trick for me. Note the need to keep any \label that may be used right
% after \section in the above as the hack puts \section within a raised box.



% The very first letter is a 2 line initial drop letter followed
% by the rest of the first word in caps (small caps for compsoc).
% 
% form to use if the first word consists of a single letter:
% \IEEEPARstart{A}{demo} file is ....
% 
% form to use if you need the single drop letter followed by
% normal text (unknown if ever used by IEEE):
% \IEEEPARstart{A}{}demo file is ....
% 
% Some journals put the first two words in caps:
% \IEEEPARstart{T}{his demo} file is ....
% 
% Here we have the typical use of a "T" for an initial drop letter
% and "HIS" in caps to complete the first word.
%\IEEEPARstart{T}{his} demo file is intended to serve as a ``starter file''
%for IEEE Computer Society journal papers produced under \LaTeX\ using
%IEEEtran.cls version 1.8 and later.
% You must have at least 2 lines in the paragraph with the drop letter
% (should never be an issue)
%

\section{El Experimento}
\subsection{Descripci\'on}\label{sec:experimento:descripcion}
\par El experimento propuesto consiste en establecer una conexi\'on entre dos
procesos (no necesariamente corriendo en el mismo \textit{host}, aunque se
realiz\'o de dicha manera por motivos que se explican m\'as adelante) utilizando
el protocolo PTC, para luego comenzar a enviar una cantidad determinada de
informaci\'on de un proceso al otro mientras se guardan los datos pertinentes
que luego ser\'a analizada. Es decir, no habr\'a transmici\'on en simult\'aneo
entre los procesos. As\'i pues, denominamos como \emph{servidor} al proceso que
recibe la informaci\'on y env\'ia los $ack's$, y \emph{cliente} al que env\'ia
la informaci\'on y recibe los $ack's$\footnote{Por comodidad de notaci\'on,
decidimos que fuese el servidor el proceso \textit{bindeado}/unido a un socket
escuchando peticiones de conexi\'on, y que fuese el cliente el encargado de
realizar el pedido de conexi\'on}.

\subsubsection*{Sobre los Datos del Experimento}
\par Los datos que se guardan durante la experimentaci\'on, para luego ser
analizados son aquellos cambios que sufre el RTO durante la transferencia del
experimento. Cada vez que el protocolo decide actualizar el RTO, se guarda en un
archivo el RTO previo al cambio, y el nuevo RTT que el protocolo considera para
actualizarlo. Es decir, la informaci\'on guardada termina representando el RTT
real que el protocolo midi\'o mediante el seguimiento de un paquete en
particular, y el RTO que se estaba utilizando para estimar dicho RTT. Como
informaci\'on adicional, tambi\'en se guarda la cantidad de paquetes env\'iados
por el cliente, cuales de ellos son retransmiciones, y cu\'antos paquetes fueron
\emph{dropeados}\footnote{Utilizamos aqu\'i la denominaci\'on paquete en lugar
de ''informaci\'on'' ya que la cantidad de paquetes que se env\'ien durante la
experimentaci\'on variar\'a seg\'un los valores de $\delta$ y $\phi$, pues PTC
trat\'a de enviar la mayor cantidad de informaci\'on posible por paquete.}
debido a los par\'ametros de simulaci\'on de congesti\'on ($\delta$ y $\phi$).

\par La informaci\'on que es transmitida desde el cliente al servidor no es de
importancia en el experimento propuesto, ya que como se puede observar, los
datos que ser\'an analizados no tienen relaci\'on alguna con el \emph{payload}
de los paquetes PTC ni con la totalidad de la informaci\'on. As\'i pues, se
decidi\'o enviar una cantidad finita de bloques de informaci\'on aleatoria de un
tama\~no determinado (tanto la cantidad de bloques como el tama\~no de los
mismos se encuentran definidos en la secci\'on
\ref{sec:variables_metricas:variables}). Con esto se obtiene un intercambio de
informaci\'on suficiente para obtener los datos del RTO y RTT necesarios para
realizar el an\'alisis posterior.

\subsubsection*{Simulaci\'on de la congesti\'on/Modificaciones en el Protocolo}
\par Ya se han mencionado en este trabajo a los par\'ametros que se utilizan
para simular la congesti\'on de la red: $\delta$ (delay) y $\phi$ (probabilidad
de p\'erdida de un paquete, o también llamada, \emph{dropchance}).

\par Si bien es cierto que conceptualmente estas caracter\'isticas de la red no
son inherentes a la congesti\'on que pueda tener una red, consideramos que son
dos de los factores m\'as importantes que nos se\~nalan que una red comienza a
estar congestionada. De hecho, TCP considera que hay congesti\'on cuando
comienza a dejar de recibir los $ack$ de sus paquetes, y el $timeout$ para
considerar que un $ack$ se ha perdido debido a la congesti\'on se calcula
din\'amicamente de acuerdo al RTT de la red (en el cual entra el juego el
delay). Entonces, obsevamos como evento que al haber congesti\'on, los paquetes
(no s\'olo los $ack's$ comienzan a ''perderse'', o mejor dicho, ser
\textit{dropeados}) y al mismo tiempo la red suele comenzar a tener un mayor
delay debido a que los distintos nodos que hay en una ruta entre dos puntos de
una conexi\'on podr\'ian estar saturados y trabajando m\'as lento, as\'i como
tambi\'en el acceso al medio podr\'ia ser m\'as lento debido a la gran cantidad
de paquetes que circulan por la red.

\par A la hora de simular estas caracter\'isticas, hay varios lugares en el
c\'odigo del protocolo donde se podr\'ia hacer. Una posibilidad ser\'ia hacer
que los $ack's$ tuvieran una demora $\delta$ para ser env\'iados, y justo antes
de enviarlos, decidir si se env\'ia o no con probabilidad $\phi$ (todo esto, en
nuestro experimento ser\'ia algo que realizar\'ia el servidor). Ahora bien, el
protocolo en muchas ocasiones tambi\'en env\'ia \emph{spurious}\footnote{La
palabra es un abuso de notaci\'on, ya que estos $ack's$ no son ni aleatorios ni
no esperados, sino que corresponden a la especificaci\'on del protocolo.} $ack$
para actualizar el tama\~no de la ventana del receptor, y la pregunta en este
caso pasar\'ia a ser: ¿debemos aplicar $\delta$ y $\phi$ para estos $ack's$
tambi\'en? ¿y c\'omo afectan a los datos que se miden durante el experimento?.

\par Originalmente, se trat\'o de trabajar con $\delta$ y $\phi$ de la manera
expuesta, pero surgieron varias complicaciones a la hora de interpretar el
\emph{sampled\_rtt}\footnote{Lo que para el protocolo es el RTT real de la red
para una conexi\'on dada.}, ya que estos paquetes de actualizaci\'on de tama\~no
de ventana pod\'ian ser utilizados por el protocolo para mandar informaci\'on de
un paquete que hab\'ia sido dropeado anteriormente. As\'i pues, en lugar de
tratar de resolver este problema y seguir con este modelo para simular el delay
y el dropchance, se decidi\'o adoptar otro enfoque.

\par En si mismo, lo que importa es simular la congesti\'on, y esta no pertenece
a una de las puntas de una conexi\'on (ni siquiera a ambas), sino que es un
evento que ocurre en la red. As\'i pues, si la misma est\'a simulada del lado
del servidor o del lado del cliente, dentro del alcance de nuestro experimento,
es exactamente lo mismo simularla en cualquiera de los dos lados.  La diferencia
radica en que el servidor tambi\'en tendr\'ia un delay a la hora de enviar
informaci\'on al cliente pudiendo afectar de alguna manera a nuestras
mediciones. Pero en el experimento planteado, el cliente no env\'ia
$ack's$\footnote{Exceptuando los correspondientes a la creaci\'on de la
conexi\'on} ni el servidor env\'ia paquetes ''regulares'', as\'i pues este
enfoque nos quita el problema de lidiar con los paquetes de actualizaci\'on de
ventana de recepci\'on.

\par Finalmente, se tom\'o el archivo \emph{protocol.py}, y se a\~nadi\'o una
probabilidad de dropeo y el delay en la funci\'on \emph{send\_and\_queue}, pero
teniendo en cuenta que todos los c\'alculos correspondientes al RTO y RTT fueran
efectuados previamente, de manera de que tuvieran el impacto que se pretende
estudiar cuando alg\'un paquete es expuesto a un delay muy grande o es
descartado/dropeado.

%-------------------------------------------------------------------------------

\subsection{Entorno de Experimentaci\'on}\label{sec:experimento:entorno}
\par El entorno donde se realizar\'on todos los experimentos analizados en este
trabajo est\'a compuesto de 6 m\'aquinas virtuales corriendo el sistema
operativo \textit{CentOS 6.5 64-bit}\cite{centos} con 2GB de memoria RAM y 2
n\'ucleos x86\_64 intel de 2.26 GHz cada una, todas corriendo sobre el mismo
host f\'isico utilizando la tecnolog\'ia de virtualizaci\'on Xen4\cite{xen}.

\par En cada uno de las VMs\footnote{M\'aquinas virtuales.} se corri\'o el
experimento descripto para una combinaci\'on distinta de valores de $\delta$ y
$\phi$, aunque para cada combinaci\'on se experiment\'o con todas las
combinaci\'on posibles para los valores de $\alpha$ y $\beta$ definidos en la
secci\'on \ref{sec:variables_metricas:variables}.

%-------------------------------------------------------------------------------

\subsection{Dificultades encontradas durante la Experimentaci\'on}\label{sec:experimento:dificultades}
\par Durante la realizaci\'on de este trabajo, varias complicaciones
relacionadas con el protocolo fueron surgiendo. Como consecuencia, algunos de
los resultados obtenidos, o mismo los experimentos realizados, no fueron tan
satisfactorios como se esperaba en cuanto a su precisi\'on.

\par A continuaci\'on, detallamos los casos m\'as importantes y las decisiones
que se tomaron en consecuencia:

\begin{enumerate}[label=(\alph*)]

    \item \textbf{Errorno 90:} Inicialmente, al comenzar con la
        experimentaci\'on, la idea era correr el servidor en un un host, y el
        cliente en otro (ambos siendo VMs). Pero al comenzar la
        experimentaci\'on nos encontramos con el error \emph{''[Errno 90]
        Message too long''}. El mismo se debe a que el protocolo, al armar los
        paquetes, no toma en cuenta\footnote{O no toma en cuenta \emph{del
        todo}.} el MTU de la interfaz de red. Este error no ocurre cuando se
        corre el servidor en la IP 127.0.0.1, ya que esta es una interfaz
        virtual que no esta asociada a ning\'un dispositivo de red, y
        aparentemente el sistema operativo no le impone un MTU.

        \par Se observ\'o que este error pod\'ia ser \emph{by-paseado}
        disminuyendo la constante \emph{MSS} a un valor por debajo del MTU del
        dispositivo de red a utilizar, pero finalmente fue otro camino el que se
        tom\'o para superar este inconveniente. Luego de hablar con los docentes
        mediante la lista de la c\'atedra, se descubri\'o que este error fue
        resuelto en uno de los branches del proyecto del protocolo
        PTC\footnote{\url{https://github.com/lukius/ptc/issues/7}}, pero al
        momento de utilizar dicho c\'odigo, el tiempo para realizar la
        experimentaci\'on de este trabajo disminu\'ia alarmantemente, por lo
        cual se decidi\'o realizar toda la experimentaci\'on en forma local (es
        decir, servidor y cliente corriendo en el mismo host) y as\'i aprovechar
        mejor las 6 VMs disponibles.

    \bigskip
    \item \textbf{Discretizaci\'on del Tiempo:} Una de las caracter\'isticas de
        PTC es que fue inspirado en la implementaci\'on de TCP de BSD\cite{bsd},
        el cual en lugar de utiliza un timer discreto en lugar del timer de alta
        precisi\'on que podr\'ia proveer el sistema operativo. En el caso de
        PTC, este timer tiene una precisi\'on de 10ms, con lo cual, si alg\'un
        paquete no es procesado por completo instant\'aneamente, el
        \emph{sampled\_rtt} correspondiente tendr\'a una cota inferior de 10ms.
        Estos 10ms que se a\~nade en algunos paquetes no es problem\'atico
        cuando
        \begin{enumerate*}[label=\itshape\arabic*\upshape)]

            \item ocurre en todos los paquetes, o

            \item cuando los par\'ametros $\delta$ y $\phi$ son bajos, con lo
                cual se esperar\'ia tener un \emph{sampled\_rtt} muy cercano a
                $0$, pero donde se obtiene o bien $0$, o $10$.

        \end{enumerate*}
        \par Para solucionar este ''inconveniente'', se podr\'ia
        rabajar con el reloj de mayor precisi\'on del sistema operativo,
        para lo cual hubiera sido necesario modificar el protocolo o implementar
        alg\'un mecanismo para utilizarlo\footnote{Utilizando los \emph{payloads
        de los paquetes quiz\'as}.}. Pero se consider\'o que estos cambios no
        terminari\'an siendo necesarios, ya que en los casos de mayor inter\'es
        (aquellos donde $\delta$ y $\phi$ tienen mayor relevancia), estos 10ms
        de m\'as que pueden aparecer en algunos paquetes tendr\'an muy poca
        influencia en el an\'alisis global.

    \bigskip
    \item \textbf{Duplicaci\'on de $\delta$:} Uno de los comportamientos que se
        observaron durante la experimentaci\'on fue que los RTT obtenidos por el
        protocolo eran, excepto para los primeros y \'ultimos paquetes del
        experimento, el doble de $\delta$, lo cual no era muy razonable ya que
        al correr todo el experimento localmente, el delay agregado por la
        ''red'' es demasiado peque\~no para influir de esta manera\footnote{De
        hecho, es de 20$\mu s$.}. Analizando el problema con
        Lucio\footnote{\url{https://github.com/lukius}}, se
        lleg\'o\footnote{\'El lleg\'o.} a la conclusi\'on que este
        comportamiento se deb\'ia un \emph{lock} que piden en simult\'aneo
        varios \emph{threads}. El recurso en cuesti\'on que se corresponde con
        el lock es el \emph{control block}, que entre otras cosas, es el que
        administra el acceso a los \emph{buffers} del socket. As\'i pues,
        tenemos al cliente que se encuentra enviando paquetes y recibiendo
        $ack's$ en simult\'aneo, y ambos procesos piden el recurso para poder
        tomar o escribir en los buffers. Entonces, al tener a $\delta$
        implementado del lado del cliente\footnote{En realidad est\'a
        implementado en el protocolo, pero a fines pr\'acticos, la funci\'on
        modificada es s\'olo invocada por el cliente en nuestro experimento.},
        se tiene el delay al enviar el paquete, y se le suma el delay del
        paquete que se este enviando (el que tenga el recurso) al momento de
        recibir el $ack$, ya que se debe esperar a que dicho paquete se env\'ie
        antes de poder obtener el \emph{control block}.

        \par Sobre este caso, esta disertaci\'on se di\'o demasiado cerca de la
        fecha de
        entrega\footnote{\url{http://www.memecreator.org/static/images/memes/2739987.jpg}},
        con lo cual se decidi\'o documentar dicho comportamiento pero procesar
        los datos tal como se obtuvieron de la experimentaci\'on. Esto lleva a
        que, en lo que sigue de este trabajo, cuando hablemos de un delay
        $\delta$, en la realidad este represent\'o un delay de $2\delta$.

    \bigskip 
    \item \textbf{Pseudo-deadlocks:} El problema m\'as grande que se enfrent\'o
        a la hora de realizar la experimentaci\'on fueron ciertos
        \emph{deadlocks} o \emph{cuelgues}. M\'as que cuelgue, ser\'ia m\'as
        preciso decir que el protocolo pasa por ciertas etapas de
        lentitud\footnote{Observamos casos de hasta 15 horas.} sin responder
        para algunas ejecuciones, independientemente de los valores de los
        par\'ametros $\alpha$, $\beta$, $\delta$, $\phi$ y el puerto. De hecho,
        se observ\'o como dos experimentos (con exactamente los mismos
        par\'ametros excepto por el puerto) ejecutados en el mismo entorno y de
        forma contigua, pueden resultar en una ejecuci\'on entrando en
        \emph{pseudo-deadlocks} y la otra finalizando correctamente en menos de
        30 segundos. Estos casos al ser independiente de los par\'ametros,
        fueron el principal impedimento para lograr una experimentaci\'on m\'as
        rica, ya que si bien son independientes, se le suma que con un
        dropchance muy alto es d\'ificil distinguir cuando el experimento
        entr\'o en este estado de $deadlock$ o si sigue trabajando normalmente.

        \par Se presume que si se dejase correr a los procesos el tiempo
        suficiente, terminar\'ian, pero esto hace que el tiempo del experimento
        sea anormalmente grande, afectando la m\'etrica del $throughput$ (que se
        introduce en la secci\'on \ref{sec:variables_metricas:metricas}).

        \par As\'i pues, se tomaron dos decisiones para poder realizar la
        experimentaci\'on. En primer lugar se decidi\'o no trabajar con
        dropchance mayor a 10\%, ya que a partir de dicho valor cada
        experimentaci\'on demanda un tiempo irrazonablemente largo\footnote{Y ya
        est\'abamos al horno con el deadline\ldots}. Y en segundo lugar, la
        experimentaci\'on se realiz\'o mediante un proceso que verifica que el
        RTO sea actualizado al menos cada 5 minutos. Pasados 5 min sin actualizar
        el RTO, se considera que el experimento entr\'o en
        \emph{pseudo-deadlock} y se cancela para luego realizarla desde cero.

\end{enumerate}
%-------------------------------------------------------------------------------


\section{Las Variables y M\'etricas}
\subsection{Valores de las variables utilizadas para la Experimentaci\'on}\label{sec:variables_metricas:variables}
\par Hasta ahora se han mencionado muchos par\'ametros para los cuales no se han
definido sus valores con los cuales luego se experimentar\'a. A continuaci\'on,
se trata de hacer res\'umen de todos estos par\'ametros y variables con sus
respectivos valores con los cuales se realiz\'o toda la experimentaci\'on:

\begin{description}

    \bigskip
    \item[$\alpha$/$\beta$:] Son los valores que utiliza el protocolo para
        calcular el RTO. Se experiment\'o, para ambos par\'ametros, con los
        valores 0, 0.1, 0.2, 0.3, 0.4, 0.5, 0.6, 0.7, 0.8, 0.9 y 1.

    \bigskip
    \item[DropChance ($\phi$):] Como ya se explic\'o, nos vimos limitados a la
        hora de variar este par\'ametro debido al tiempo requerido por cada
        experimento con un dropchance por encima del 10\%. As\'i pues, se
        logr\'o experimentar con 0\%, 1\%, 5\% y 10\%.

    \bigskip
    \item[Delay ($\delta$):] El delay inyectado a cada experimento trat\'o no
        solo de simular el delay que puede producirse mediante la congesti\'on,
        sino tambi\'en aquellas redes de mayor extensi\'on o que debidos a sus
        medios de propagaci\'on se ven afectados de manera de tener un mayor
        delay. Se ensay\'o con los valores: 0ms, 5ms, 10ms, 25ms, 50ms, 100ms,
        250ms y 500ms.

    \bigskip
    \item[Tama\~no los mensajes:] El experimento se basa en trasferir
        informaci\'on del cliente al servidor. Si bien la informaci\'on en si
        que se transmite no es de importancia, para generala se realizan
        mensajes con un tama\~no de 500bytes de informaci\'on aleatoria, los
        cuales luego son procesados por el protocolo que puede agrupar a varios
        (o partes) de estos en los paquetes PTC.

    \bigskip
    \item[Cantidad de mensajes:] As\'i como el dropchance afecta negativamente a
        la duraci\'on de cada experimento, es claro que la cantidad de
        informaci\'on que se pretende env\'iar tambi\'en lo har\'a. As\'i pues,
        para tener una muestra significativa, por cada experimento (es decir,
        para las dem\'as variables ya fijadas) se transmitieron 100
        mensajes\footnote{Originalmente eran 1000, pero esto hizo imposible la
        experimentaci\'on de probabilidades con un delay superior a 10ms, o con
        un dropchance mayor a 5\%.}.

\end{description}

%-------------------------------------------------------------------------------

\subsection{M\'etricas Utilizadas}\label{sec:variables_metricas:metricas}
\par Una vez terminada la experimentaci\'on (o previamente en realidad, para
determinar que informaci\'on es necesaria reunir), queda pendiente decidir como
procesar la misma para poder visualizarla y analizarla de alguna manera. Para
ello, se utilizaron las siguientes m\'etricas: 

\begin{description}
    \bigskip
    \item[\textbf{Root Mean Square Deviation:}] El foco del enunciado est\'a
        puesto en las variables que se utilizan para calcular el RTO (es decir,
        la estimaci\'on del RTT del protocolo). As\'i pues, concentrados en cuan
        ''bien'' estima el protocolo el RTT, nos es interesante encontrar alguna
        manera de poder comparar cada estimaci\'on con el que luego fue el RTT
        real correspondiente. Ahora bien, esto se podr\'ia realizar para cada
        estimaci\'on realizada por el protocolo (que si bien no son para todos
        los paquetes que env\'ia, son muchos), hay que encontrar una forma de
        unificar estos valores para decidir, en alg\'un tipo de promedio, cuan
        bueno es el protocolo con ciertos par\'ametros estimando el RTT.

        \par Teniendo tanto los valores de RTO estimados durante la
        experimentaci\'on, como as\'i tambi\'en el RTT que se correspondi\'o con
        cada RTO (es decir, el valor que se estaba estimando), buscamos alguna
        m\'etrica de distancia promedio entre una serie de valores estimados y
        otra de valores ''ideales'' o reales. As\'i llegamos al
        \emph{RMSD}\footnote{Y Rodrigo no tuvo nada que ver con
        esto.}\cite{rmsd}, que justamente es una de las m\'etricas conocidas
        para este tipo de ensayos. Mientras m\'as cercano a cero sea esta
        m\'etrica, mejor es la estimaci\'on, ya que esta es un valor
        estandarizado del desv\'io entre la estimaci\'on y el valor real.

    \bigskip
    \item[\textbf{Throughput:}] El \emph{throughput}\cite{hroughput} es
        b\'asicamente una m\'etrica de rendimiento, que normalmente en el \'area
        de las redes y comunicaciones, refiere a la cantidad de informaci\'on
        efectiva transmitida sin errores por unidad de tiempo.
        En nuestro caso, conociendo la cantidad de mensajes que se
        transmitir\'an, y el tama\~no de los mismos, ya sabemos de antemano
        cuanta informaci\'on efectiva (medida en bytes), es decir informaci\'on
        que quiere transmitir la aplicaci\'on que utilice el protocolo, se
        transmitir\'an en todo el experimento. As\'i pues, sabiendo cuanto
        tiempo llev\'o completar el experimento para un conjunto de
        par\'ametros, se puede saber cual fue el througput obtenido, el cual
        refleja la percepci\'on real de la aplicaci\'on sobre cuan bueno es el
        protocolo.

    \bigskip
    \item[\textbf{Compromiso RMSD-Throughput:}] Por \'ultimo, se desea observar
        la relaci\'on de compromiso entre la estimaci\'on del RTO y el
        throughput, para cada conjunto dado. Podr\'ia darse el caso que quiz\'as
        una combinaci\'on particular de $\alpha$ y $\beta$ para alg\'un tipo de
        red (con mucho delay y dropchance, por ejemplo) no obtuviera el mejor
        throughput ni el mejor RMSD, pero si obtuviese el mejor balance entre
        estos dos.

        \par Siendo el RMSD un valor que es mejor en la medida que m\'as cerca
        del cero este, y el throughput una medida que mientras m\'as grande sea
        nos comunica un mejor rendimiento, es claro que observar la m\'etrica
        $Throughput / RMSD$ nos dar\'a un resultado del compromiso entre estos
        dos objetivos deseados. Mientras m\'as alto sea esta m\'etrica, mejor
        ser\'a la combinaci\'on de par\'ametros deseados a la hora de buscar
        maximizar ambas caracter\'isticas de forma balanceada.

\end{description}

%-------------------------------------------------------------------------------


\section{Resultados y An\'alisis}
\subsection{RMSD}\label{sec:resultados:rmsd}
\newpage
\foreach \prob in {0, 1, 5, 10}{
    \foreach \delay in {0,5,10,25,50,100,250,500}{
        \begin{figure}
            \centering
            \includegraphics[angle=-90,width=0.5\textwidth]{img/rmsd/prob\prob_delay\delay_rmsd.eps}
            \caption{Saraza}
        \end{figure}
    }
}

%-------------------------------------------------------------------------------

\subsection{Throughput}\label{sec:resultados:throughput}

%-------------------------------------------------------------------------------

\subsection{Compromiso RMSD-Throughput}\label{sec:resultados:rms_vs_throughput}

%-------------------------------------------------------------------------------


%\subsection{Subsection Heading Here}
%Subsection text here.

% needed in second column of first page if using \IEEEpubid
%\IEEEpubidadjcol

%\subsubsection{Subsubsection Heading Here}
%Subsubsection text here.


% An example of a floating figure using the graphicx package.
% Note that \label must occur AFTER (or within) \caption.
% For figures, \caption should occur after the \includegraphics.
% Note that IEEEtran v1.7 and later has special internal code that
% is designed to preserve the operation of \label within \caption
% even when the captionsoff option is in effect. However, because
% of issues like this, it may be the safest practice to put all your
% \label just after \caption rather than within \caption{}.
%
% Reminder: the "draftcls" or "draftclsnofoot", not "draft", class
% option should be used if it is desired that the figures are to be
% displayed while in draft mode.
%
%\begin{figure}[!t]
%\centering
%\includegraphics[width=2.5in]{myfigure}
% where an .eps filename suffix will be assumed under latex, 
% and a .pdf suffix will be assumed for pdflatex; or what has been declared
% via \DeclareGraphicsExtensions.
%\caption{Simulation Results.}
%\label{fig_sim}
%\end{figure}

% Note that IEEE typically puts floats only at the top, even when this
% results in a large percentage of a column being occupied by floats.
% However, the Computer Society has been known to put floats at the bottom.


% An example of a double column floating figure using two subfigures.
% (The subfig.sty package must be loaded for this to work.)
% The subfigure \label commands are set within each subfloat command,
% and the \label for the overall figure must come after \caption.
% \hfil is used as a separator to get equal spacing.
% Watch out that the combined width of all the subfigures on a 
% line do not exceed the text width or a line break will occur.
%
%\begin{figure*}[!t]
%\centering
%\subfloat[Case I]{\includegraphics[width=2.5in]{box}%
%\label{fig_first_case}}
%\hfil
%\subfloat[Case II]{\includegraphics[width=2.5in]{box}%
%\label{fig_second_case}}
%\caption{Simulation results.}
%\label{fig_sim}
%\end{figure*}
%
% Note that often IEEE papers with subfigures do not employ subfigure
% captions (using the optional argument to \subfloat[]), but instead will
% reference/describe all of them (a), (b), etc., within the main caption.


% An example of a floating table. Note that, for IEEE style tables, the 
% \caption command should come BEFORE the table. Table text will default to
% \footnotesize as IEEE normally uses this smaller font for tables.
% The \label must come after \caption as always.
%
%\begin{table}[!t]
%% increase table row spacing, adjust to taste
%\renewcommand{\arraystretch}{1.3}
% if using array.sty, it might be a good idea to tweak the value of
% \extrarowheight as needed to properly center the text within the cells
%\caption{An Example of a Table}
%\label{table_example}
%\centering
%% Some packages, such as MDW tools, offer better commands for making tables
%% than the plain LaTeX2e tabular which is used here.
%\begin{tabular}{|c||c|}
%\hline
%One & Two\\
%\hline
%Three & Four\\
%\hline
%\end{tabular}
%\end{table}


% Note that IEEE does not put floats in the very first column - or typically
% anywhere on the first page for that matter. Also, in-text middle ("here")
% positioning is not used. Most IEEE journals use top floats exclusively.
% However, Computer Society journals sometimes do use bottom floats - bear
% this in mind when choosing appropriate optional arguments for the
% figure/table environments.
% Note that, LaTeX2e, unlike IEEE journals, places footnotes above bottom
% floats. This can be corrected via the \fnbelowfloat command of the
% stfloats package.



%\section{Conclusion}
%The conclusion goes here.

\section{Conclusiones}\label{sec:conclusiones}
    \IEEEPARstart{LL}egados al final del trabajo expuesto, podemos decir que determinar valores 
de $\alpha$ y $\beta$ que acerquen la estimac\'ion del RTO al RTT ''real'' no es tarea sencilla.

En base a los resultados obtenidos mediante la primera m\'etrica que usamos, el RMSD, tener 
tantos valores variables en consideraci\'on nos genera un campo demasiado amplio de resultados.
A pesar de eso pudimos ver la relaci\'on directa que tienen los cambios introducidos al
protocolo (el delay y la probabilidad de dropeo) en la estimaci\'on del RTO como bien describimos
en la secci\'on de resultados.

Algo similar sucedi\'o al considerar la segunda m\'etrica de la experimentac\'on, el throughput,
a grandes rasgos no se consegu\'ia ver una relaci\'on clara entre la informaci\'on transmitida sin
errores y los valores de $\alpha$ y $\beta$ como s\'i suced\'ia con $\phi$ y $\delta$, los resultados 
obtenidos reflejaban como era de esperarse que a mayor delay menor era la eficiencia de la comunicac\'ion 
y que a mayor probabilidad de p\'erdida de paquete lo mismo suced\'ia.


Pero no fue hasta analizar enfonc\'andose en un entorno m\'as peque\~no de valores, esto es,
al fijar ciertos par\'ametros y concentrarnos en como repercut\'ian $\alpha$ y $\beta$, que pudimos 
determinar valores que consideramos eficientes de ellos dependiendo el caso.

A partir de las modificaciones hechas al protocolo, el estudio de distintas m\'etricas y formas de
presentar y analizar los resultados, se reduci\'o el tiempo que pudimos dedicarle a la experimentaci\'on
propiamente dicha, y quedaron pendientes posibles trabajos a realizar, uno de ellos es estudiar el
compromiso entre el RMSD y el throughput [AMPLIAR CON LO QUE ESTA EN PAG6.]

Otro punto que vale la pena destacar es que los resultados obtenidos dependen completamente de la 
implementaci\'on del protocolo, que consideramos no finalizada, ya que hubiese sido interesante
continuar modificando el c\'odigo a partir de los errores que llegamos a reportar en la secci\'on 
\ref{sec:experimento:dificultades}, pero no resolver completamente, lo cual nos hubiese permitido 
a su vez estudiar distintos entornos de experimentaci\'on y hacer un an\'alisis m\'as profundo de 
multiples casos posibles. Al continuar con esto, los resultados pueden llegar a variar y ser
menos vol\'atiles o irregulares que los que presentamos.

Como palabras finales sobre el trabajo hecho, consideramos que el estudio de la congesti\'on de TCP
es m\'as que interesante y el trabajo presentado nos permite ver en pr\'actica la importancia de varios 
conceptos de la capa de transporte y el control de congestion estudiados a lo largo de la cursada
de la materia.



% if have a single appendix:
%\appendix[Proof of the Zonklar Equations]
% or
%\appendix  % for no appendix heading
% do not use \section anymore after \appendix, only \section*
% is possibly needed

% use appendices with more than one appendix
% then use \section to start each appendix
% you must declare a \section before using any
% \subsection or using \label (\appendices by itself
% starts a section numbered zero.)
%


%\appendices
%\section{Proof of the First Zonklar Equation}
%Appendix one text goes here.

% you can choose not to have a title for an appendix
% if you want by leaving the argument blank
%\section{}
%Appendix two text goes here.


% use section* for acknowledgement
%\ifCLASSOPTIONcompsoc
  % The Computer Society usually uses the plural form
%  \section*{Acknowledgments}
%\else
  % regular IEEE prefers the singular form
%  \section*{Acknowledgment}
%\fi


%The authors would like to thank...


% Can use something like this to put references on a page
% by themselves when using endfloat and the captionsoff option.
\ifCLASSOPTIONcaptionsoff
  \newpage
\fi



% trigger a \newpage just before the given reference
% number - used to balance the columns on the last page
% adjust value as needed - may need to be readjusted if
% the document is modified later
%\IEEEtriggeratref{8}
% The "triggered" command can be changed if desired:
%\IEEEtriggercmd{\enlargethispage{-5in}}

% references section

% can use a bibliography generated by BibTeX as a .bbl file
% BibTeX documentation can be easily obtained at:
% http://www.ctan.org/tex-archive/biblio/bibtex/contrib/doc/
% The IEEEtran BibTeX style support page is at:
% http://www.michaelshell.org/tex/ieeetran/bibtex/
%\bibliographystyle{IEEEtran}
% argument is your BibTeX string definitions and bibliography database(s)
%\bibliography{IEEEabrv,../bib/paper}
%
% <OR> manually copy in the resultant .bbl file
% set second argument of \begin to the number of references
% (used to reserve space for the reference number labels box)
%\begin{thebibliography}{1}

%\bibitem{IEEEhowto:kopka}
%H.~Kopka and P.~W. Daly, \emph{A Guide to \LaTeX}, 3rd~ed.\hskip 1em plus
%  0.5em minus 0.4em\relax Harlow, England: Addison-Wesley, 1999.

%\end{thebibliography}

% biography section
% 
% If you have an EPS/PDF photo (graphicx package needed) extra braces are
% needed around the contents of the optional argument to biography to prevent
% the LaTeX parser from getting confused when it sees the complicated
% \includegraphics command within an optional argument. (You could create
% your own custom macro containing the \includegraphics command to make things
% simpler here.)
%\begin{IEEEbiography}[{\includegraphics[width=1in,height=1.25in,clip,keepaspectratio]{mshell}}]{Michael Shell}
% or if you just want to reserve a space for a photo:

%\begin{IEEEbiography}{Michael Shell}
%Biography text here.
%\end{IEEEbiography}

% if you will not have a photo at all:
%\begin{IEEEbiographynophoto}{John Doe}
%Biography text here.
%\end{IEEEbiographynophoto}

% insert where needed to balance the two columns on the last page with
% biographies
%\newpage

%\begin{IEEEbiographynophoto}{Jane Doe}
%Biography text here.
%\end{IEEEbiographynophoto}

% You can push biographies down or up by placing
% a \vfill before or after them. The appropriate
% use of \vfill depends on what kind of text is
% on the last page and whether or not the columns
% are being equalized.

%\vfill

% Can be used to pull up biographies so that the bottom of the last one
% is flush with the other column.
%\enlargethispage{-5in}



% that's all folks
\onecolumn
\printbibliography
\end{document}
